\section{Example: APPLE-II Undulator} \label{ex4}

The most realistic undulator model
\footnote{See also example \ref{ex11} ~ for an even better model by interfacing
to \Ue. \U allows to take into account the interaction
between the neighbouring magnets, while WAVE only consideres static
magnets with $\mu = 1$. This interaction might affect especially the
first and second integrals, i.e. the kick and off set at the end of
the trajectory.}
of \W is an undulator build of
permanent magnet blocks (rare earth componds like NdFeB, shortly REC).
The field of these blocks is calculated with the current sheet method.
Due to the large number of blocks of an undulator the field
calculations are rather slow, but the field is very realistic, and an
undulator can be modelled with all its features like gap and shift
changes.

The namelist \scap{\$REC} holds the variables of the undulator. It also
includes gap taper and field errors. To set up the input file \wie, enter:

\cboxs{
\senv\\ {\grau
python3 ../python/wave\_example\_4.py\\
cp \dw/input/wave\_example\_4.in wave.in}
}

After the \W run, you'll find a file named \gtis{rec\_plotm.eps}, which
shows the set-up of the undulator (Fig. \ref{ex4fig1}). \\[-2cm]

\begin{figure}[ht]
    \centering
    \includegraphics[width=0.8\textwidth]
    {/home/scheer/wav/doc/pictures/wave_example_4_recplotm_1.pdf}
    \vspace{-1.5cm}
    \caption{REC model of an undulator}
    \label{ex4fig1}
\end{figure}

The figur shows the upper and lower rows of magnets as top views and a
side view of both. To see details we
click \gtis{Conf. Mag.} in the toolbar of \WSH and \\

\witem{REC-undulator (Option KBREC)}
{~}\\[-1.5cm]
\witem{Options}
{~}\\[-1.5cm]
\witem{Plotting options}

and set the plotting range

\cboxss{RPLXMN = -920.\\
RPLXMX = -850.
}

and rerun \W.\\

\begin{figure}[ht]
{~}\\[-2.5cm]
    \centering
    \includegraphics[width=0.8\textwidth]
    {/home/scheer/wav/doc/pictures/wave_example_4_recplotm_2.pdf}
    \vspace{-1.5cm}
    \caption{Endpoles of the REC model}
    \label{ex4fig2}
\end{figure}

Now, you see the details of the endpoles, the shift, and the
magnetization of the blocks. More complex situtation can be handle by
an input file containing arrays of magnet. This will be explained in
one of the following examples.

The overview of this run is shown in Fig. \ref{ex4fig3}. It shows the
elliptical trajectory and the corresponding first harmonic.\\

\begin{figure}[ht]
    \centering
    %\framebox(11.5,8){
      \includegraphics[width=0.7\textwidth]
      {/home/scheer/wav/doc/pp/wave_example_4_overview.pdf}
    %}
    \caption{Example of an APPLE-II undulator}
    \label{ex4fig3}
\end{figure}

\newpage
\underline{Proposed exercise:}

In the option menu of the REC-Undulators is a menu to define errors of
the magnetic blocks. Give it a try, and check also the magnetic field
errors option in the menu of the global magnetic field options.
